\begin{tikzpicture}[shorten >=1pt,<->,draw=black!25, node distance=\layersep]
    \tikzstyle{every pin edge}=[<-,shorten <=1pt]
    \tikzstyle{neuron}=[circle,fill=black!25,minimum size=12pt,inner sep=0pt]
    \tikzstyle{input neuron}=[neuron, fill=blue!50];
    \tikzstyle{hidden neuron}=[neuron, fill=red!50];
    \tikzstyle{annot} = [text width=4em, text centered]

    % Draw the input layer nodes
    \foreach \name / \y in {1,...,5}
    % This is the same as writing \foreach \name / \y in {1/1,2/2,3/3,4/4}
    \node[input neuron] (I-\name) at (0,-\y cm) {};

    % Draw the hidden layer nodes
    \foreach \name / \y in {2,...,4}
   \node[hidden neuron] (H-\name) at (\layersep,-\y cm) {};

    % Connect every node in the input layer with every node in the
    % hidden layer.
    \foreach \source in {1,...,5}
        \foreach \dest in {2,...,4}
        \path (I-\source) edge (H-\dest);

    % Connect every node in the input layer with node 3 in the
    % hidden layer in black
    \foreach \source in {1,...,5}
        \path (I-\source) edge[black] (H-3);
        
    % Annotate the layers
    \node[annot,above of=I-1, node distance=1cm] (vl) {Visual layer};
    \node[annot,right of=vl] {Hidden layer};
\end{tikzpicture}
